% Section 7: general conclusion across the three methods.
% No new measurements; the numbers stay in Sections 2 to 5.

\section{Conclusion}

The three methods left us with three different results. FOMM moved head turns
and blinks smoothly and produced its clip in about 2.5 minutes, but it lost rare
details such as braces when it inpainted the mouth. The rectified flow pipeline
produced a 512 $\times$ 512 clip with sound on the 4 GB laptop we had, and its
restoration and smoothing stages removed much of the boil. Our own autoencoder
rebuilt faces well, but it never learned to move an identity from one person to
the other.

The failures had reasons, or at least clear hypotheses. The autoencoder loss
went down while the swap stayed wrong, because the training objective rewarded
reconstruction and our dataset of 238 images was small; 100 more epochs did not
change the result. The lost braces on FOMM have a hypothesis and no direct
evidence: the training videos may show no braces, so the learned appearance
priors cannot rebuild them. In the rectified flow clip every frame is generated
on its own, so the face boils, and the smoothing only works inside short
windows; the worst frames still jump. The newer adapter v1.1 was twelve times
slower than v1, so the recorded clip kept v1.

What we learned reaches past the single results. A low loss does not mean a
usable face swap, so the measure must match the goal. A model built for the
task, even a released checkpoint, gave more than our own simple model. Per-frame
generation always boils, and only a pass across frames reduces it. Finally,
every measure here has limits: the face change value mixes motion, lighting,
texture and flicker, and it has no calibrated threshold, so we read the numbers
together with the frames, not instead of them.
